\section{Chương~\ref{sec:pain}. Đau}

Cảm biến đau, hai dây, cổng trong tủy sống, núm âm lượng đi xuống, nhạy cảm hóa và việc chiếm sự chú ý đã được đo trực tiếp; các công thức cổng và nhạy cảm hóa là mô hình đơn giản hóa của sách; quy tắc mà cỗ máy dùng để chọn âm lượng báo động là giả thuyết.

{\footnotesize
\setlength{\tabcolsep}{3pt}\renewcommand{\arraystretch}{1.15}
\begin{longtable}{>{\raggedright}p{0.5cm}>{\raggedright}p{4.0cm}>{\raggedright}p{5.6cm}>{\raggedright}p{2.3cm}>{\raggedright\arraybackslash}p{1.7cm}}
\toprule
TT & Khẳng định & Thí nghiệm và điều đã đo & Nguồn & Trạng thái \\
\midrule\endfirsthead
\toprule
TT & Khẳng định & Thí nghiệm và điều đã đo & Nguồn & Trạng thái \\
\midrule\endhead
1 & Ngưỡng cao: ngưỡng nhiệt của các cảm biến đau khác nhau từ $41^\circ$ đến $46$--$48^\circ$ & Ghi từng sợi thần kinh. Da khỉ: ngưỡng trung vị của cảm biến chậm (C) là $41^\circ$, của cảm biến nhanh loại hai là $46^\circ$, loại một trên $53^\circ$. Da người: cảm biến C đáp ứng cả với áp lực là $40.7^\circ$, không đáp ứng với áp lực là $48^\circ$. & \cite{Treede1995heat, Weidner1999cfib} & đã đo \\
2 & Cảm biến đau thích nghi yếu hơn nhiều so với cảm biến xúc giác và sau tổn thương nhẹ trở nên nhạy hơn; suy giảm sau khi bị nung nóng mạnh đã đo ở một loại & Bỏng da nhẹ ($50^\circ$, $100\,\text{s}$): sau $5$--$10$~phút, ngưỡng đau ở người và ngưỡng của cảm biến C ở khỉ giảm $2$--$6^\circ$; các cảm biến da khác không có sự dịch chuyển như vậy. Ở cảm biến nhanh loại hai, đáp ứng sau khi nung nóng mạnh bị ức chế. Việc so sánh sự thích nghi với cảm biến xúc giác là một ước tính chung, ở đây không có thí nghiệm riêng hỗ trợ. & \cite{LaMotte1982hyper, Treede1995heat} & đã đo \\
3 & Hai dây: nhanh $5$--$30$~m/s, chậm $0.5$--$2$~m/s; thời gian đến $t=L/v$~\eqref{eq:paintime} & Tốc độ dẫn truyền: cảm biến đau nhanh của khỉ trung bình $15$ và $25$~m/s (hai loại); cảm biến C của người $0.8$--$1.0$~m/s. Với $L=1$~m, đó là vài chục mili giây và khoảng một giây. & \cite{Treede1995heat, Weidner1999cfib} & đã đo \\
4 & Cơn đau đến hai lần: trước là đau nhanh và chính xác, sau là đau âm ỉ và kéo dài & Thời gian phản ứng với nhiệt ở cẳng tay người: từ $1100$ xuống $400\ms$ khi nhiệt độ tăng; nung nóng mạnh da có lông cho đau kép, lòng bàn tay thì không. Cơn đau thứ nhất chỉ có thể do các sợi nhanh hơn $6$~m/s mang; theo thời gian tiềm, ứng với nó là cảm biến nhanh loại hai, loại không có ở lòng bàn tay. Việc phân vai (``ở đâu'' và ``tệ đến mức nào'') là cách diễn giải. & \cite{CampbellLaMotte1983first, Treede1995heat} & đã đo \\
5 & Cổng: nơ-ron đầu ra của tủy sống nhận đau và chạm, nơ-ron trung gian ức chế được cái chạm kích thích (hình~\ref{fig:gate}) & Sơ đồ cổng được đề xuất như một lý thuyết. Chuột nhắt, đánh dấu di truyền và loại bỏ các nhóm nơ-ron: nơ-ron kích thích chứa somatostatin cần cho đau cơ học; nơ-ron ức chế chứa dynorphin ngăn đầu vào từ cảm biến xúc giác tới được chúng; không có các nơ-ron này, cái chạm nhẹ gây đau. & \cite{MelzackWall1965, Duan2014} & đã đo \\
6 & Cái chạm làm dịu cơn đau & Người, xung laser gây đau trên cánh tay: cái chạm đồng thời làm giảm mức đau được đánh giá và khả năng phân biệt năng lượng xung; hiệu ứng yếu đi theo khoảng cách giữa điểm chạm và điểm đau trong cùng một đoạn da. & \cite{Mancini2014touch} & đã đo \\
7 & Giả định của thuyết cổng: đau mạnh tự tắt nơ-ron trung gian ức chế, và cổng mở toang & Trong chương được ghi rõ là giả định của thuyết ban đầu. Công trình trên chuột nhắt mô tả cách cổng chặn cái chạm đi vào kênh đau; việc cơn đau tắt nơ-ron trung gian không được chứng minh trong đó. & \cite{MelzackWall1965} & giả thuyết \\
8 & Cổng~\eqref{eq:gate}: ngưỡng $0.18$ khi không chạm, $0.86$ khi có chạm, $0.11$ khi $g=2$; khi $\nu=1$, cái chạm giảm đầu ra từ $1$ xuống $0.25$, khi $\nu=2$ không có tác dụng; bản thân cái chạm không đi qua (hình~\ref{fig:gatecurves}) & Tính theo \texttt{models/\allowbreak pain\_\allowbreak gate.py} với trọng số trong bài tái hiện mọi con số này. & \texttt{models/\allowbreak pain\_\allowbreak gate.py} & mô hình \\
9 & Các đường đi xuống từ thân não thay đổi hệ số khuếch đại của cổng theo cả hai chiều & Chuột cống, hành não, ghi khi nung nóng đuôi: một số nơ-ron phóng điện trước khi rụt đuôi (``bật''), số khác im lặng (``tắt''); kích thích yếu vùng này ức chế phản xạ. Tổng quan: cùng các mạch ấy vừa tạo thuận lợi, vừa ức chế việc truyền cơn đau, và opioid tác động thông qua chúng. & \cite{Fields1983rvm, Fields2004} & đã đo \\
10 & Sự trông đợi được đỡ đau làm dịu cơn đau qua kênh opioid & Người bị đau cấp: ở những người mà sự trông đợi làm giảm đau, việc chặn thụ thể opioid đưa cơn đau trở lại mức của những người còn lại; ở những người còn lại, việc chặn không làm đổi cơn đau. & \cite{Levine1978} & đã đo \\
11 & Não chọn âm lượng báo động theo lợi ích của việc ngắt & Chưa có thí nghiệm nào đo được quy tắc chọn âm lượng. & --- & giả thuyết \\
12 & Nhạy cảm hóa: sau tổn thương, đầu ra của cổng tăng, ngưỡng giảm, một phần là do tủy sống; nó vẫn còn sau khi kích thích gây tổn thương đã ngừng & Chuột cống: sau tổn thương da, ngưỡng phản xạ gập giảm, còn đáp ứng tăng; phân tích điện sinh lý cho thấy một phần sự tăng nảy sinh ngay trong tủy sống. Kích thích gây tổn thương gây ra rối loạn cảm giác kéo dài, tồn tại lâu hơn chính kích thích. Chương không nêu thời gian suy giảm chính xác. & \cite{Woolf1983} & đã đo \\
13 & Nhạy cảm hóa~\eqref{eq:sensitization} là phản hồi dương; khi vòng mạnh, báo động có thể tự chốt sau khi vết thương đã lành & Suy ra từ mô hình~\eqref{eq:sensitization} (bài tập cuối chương). Chưa có thí nghiệm cho thấy cơn đau dai dẳng sau khi lành là một vòng tự chốt đúng kiểu này. & suy ra trong chương & giả thuyết \\
14 & Cơn đau là phần thưởng âm, và việc học từ nó đi theo sai số sai phân thời gian & Máy chụp cộng hưởng từ, người, chuỗi tín hiệu báo trước cơn đau: hoạt động của thể vân bụng và vỏ não thùy đảo trước trùng với tín hiệu học theo sai phân thời gian. Khỉ: một phần nơ-ron \nt{DOPA} bị kích thích bởi luồng hơi khó chịu và tín hiệu báo trước nó, phần khác bị ức chế. & \cite{Seymour2004, Matsumoto2009} & đã đo \\
15 & Cơn đau ngắt công việc đang làm & Người, kích thích điện gây đau và các phép thử âm thanh: đáp ứng với phép thử $100\ms$ sau khi cơn đau bắt đầu chậm đi rõ rệt, sau $1.5\,\text{s}$ không còn khác biệt với đối chứng. Bài tổng quan tập hợp các thí nghiệm như vậy thành mô hình chiếm sự chú ý. & \cite{Crombez1996disrupt, EcclestonCrombez1999} & đã đo \\
16 & Phản xạ rụt khép kín trong tủy sống & Chuột cống: nơ-ron ở các lớp sâu của sừng sau lặp lại khuôn mẫu không gian của phản xạ rụt của từng cơ; không thể kích thích ngược chiều chúng từ đoạn cổ, tức đó là các nơ-ron trung gian tủy sống. & \cite{Schouenborg1995wr} & đã đo \\
\bottomrule
\end{longtable}}
